%% kapa: exam=1 round=26 year=2015 period=1 form=A subject=economics subject_ko=경제원론 seq=20 number=60
%% source: Q-Net 감정평가사 기출문제 (cst00309) artlSeq=5138001 file=extracted/1008/5138001/제26회_감정평가사 1차_기출/제26회_감정평가사 1차_기출/1교시_감평_A형.pdf
%% conversion: kordoc 4.24.51 -> tools/gen1.py (auto); review=proofread-sonnet; flags=none
\begin{question}{60}
중고차 시장에서 품질에 대한 정보의 비대칭성이 존재하는 경우 나타날 수 있는 현상으로 옳은 것을 모두 고른 것은?
\begin{framedtext}
ㄱ. 정보의 비대칭성이 없는 경우보다 시장에서 거래되는 중고차의 품질이 올라간다. \par
ㄴ. 보증(warranty)과 같은 신호발송(signaling)을 통해 정보의 비대칭으로 인한 문제를 완화할 수 있다. \par
ㄷ. 역선택(adverse selection)의 문제가 발생할 수 있다. \par
\end{framedtext}
\begin{choices}
  \item ㄱ
  \item ㄴ
  \item ㄱ, ㄴ
  \item ㄴ, ㄷ
  \item ㄱ, ㄴ, ㄷ
\end{choices}
\end{question}
\answer{4}
